\chapter{动态规划、MDP/POMDP 与强化学习规划}
\label{ch:dp-rl}
\chapterstatus{结构已定，核心推导待展开}

\begin{learningobjectives}
从 Bellman optimality 推导 value/policy iteration；理解 MDP 与 POMDP 的状态、动作、转移、观测和 reward；区分 model-based planning 与 model-free RL。
\end{learningobjectives}

\begin{prerequisites}
前置为第~\ref{ch:probability}、\ref{ch:path-planning}、\ref{ch:behavior}~章。实验先在简化赛场建立精确 DP baseline，再与 RL 比较样本效率、域外泛化和安全约束。基础参考\cite{sutton2018reinforcement,bertsekas2017dynamic}。
\end{prerequisites}

\plannedsection{Bellman principle}{从最短路的 cost-to-go 到有限时域动态规划。}
\plannedsection{MDP}{状态充分性、Markov 性、策略、占用分布和价值函数。}
\plannedsection{POMDP 与 belief}{观测不等于状态，belief update 与近似规划。}
\plannedsection{强化学习}{policy/value/model learning 各自替代未知的哪一部分；offline/online 数据边界。}
\plannedsection{安全与 sim-to-real}{约束、shield、domain randomization、offline validation 和人工接管。}

\begin{failuremode}
把规则计分直接写成 reward 不保证得到可部署策略。部分奖励会引入 specification gaming；训练分布覆盖不足时，高仿真回报也不能证明真实对抗中的安全和泛化。
\end{failuremode}

\begin{chaptertransition}
单机器人 MDP 忽略了队友通信和对手策略。下一章把 assignment、网络信息和博弈结构加入状态与决策，形成雷达站--哨兵--英雄等多机器人协同问题。
\end{chaptertransition}

